\section{Related Work}
We review prior work in design patterns, libraries and grammars, and interactive design tools for composite visualization.
\subsection{Design Patterns of Composite Visualization}

Composite visualizations integrate multiple visual representations to support
complex analytical and communicative goals and are widely used in visual
analytics systems~\cite{vasurvey, Survey16, deng2023survey, AI4VIS}. Compared
with single-view visualizations, their design requires consideration of not
only the encodings within individual charts but also the spatial, semantic,
and data relationships among charts~\cite{nested, lyi2020comparative,
Composite, vaid}.

Prior work has characterized composite visualizations through a range of
composition patterns. Early taxonomies identify high-level patterns such as
juxtaposition, superimposition, overloading, and nesting
~\cite{javed2012exploring}. Subsequent work provides more fine-grained
characterizations of their visual variations~\cite{Composite}, while recent
schema-based representations formalize composite designs using operators such
as concatenation, layering, faceting, and nesting and capture the data types
encoded by their components~\cite{vaid}. Together, these studies establish a
useful vocabulary for describing composite visualization structures.

However, these representations are primarily designed to characterize
existing visualizations rather than support their construction. Interactive
authoring additionally requires explicit decisions about the granularity of
reusable components and the conditions under which they can be composed. For
example, visually related chart variants may require different data
dimensions, while composition operations may impose constraints on shared
fields, coordinate systems, or data hierarchies. Leaving these requirements
abstract makes descriptive patterns difficult to translate directly into
interactive authoring operations.

{\system} builds on these taxonomies by operationalizing their composition
patterns for authoring. It defines reusable chart-level building blocks
together with their data-dimensional requirements and enables users to compose
them through direct manipulation.


\subsection{Libraries and Grammars for Composite Visualization}

Visualization libraries~\cite{Protovis, Atlas, Encodable, ProtoGraph, bostock2011d3} and grammars~\cite{nebula, Atlas, Decinteraction, satyanarayan2015reactive, satyanarayan2016vega, zong2022animated, VolumeVisualization, moreira2023urban, Atom, Gosling, dmutidens, Libra} provide formal mechanisms for authoring visualizations. 

In composite visualization authoring, low-level libraries such as D3~\cite{bostock2011d3} offer fine-grained control over different visual marks, but require authors to manually implement relationships among them, including alignment, layout coordination, and shared scales. Higher-level declarative grammars such as Vega-Lite~\cite{satyanarayan2016vega} reduce some of this implementation burden by providing operators and configuration options for common cases such as concatenation, faceting, and layering. Nevertheless, these compositions are still specified through textual or hierarchical structures, rather than through direct associations among visual objects. When compositions become complex, authors often need high familiarity with the specification structure and produce lengthy hierarchical descriptions. Hence, expressing visual composition intents through code or declarative specifications remains indirect for iterative canvas-based arrangement, nesting, layering, and refinement.

{\system} adopts a complementary approach by maintaining an internal schema for composite visualization authoring. The schema represents composition relationships among visual components and is incrementally updated through direct manipulation. In this way, users can construct complex composite structures through direct manipulation, without manually adjust complex specifications.

\subsection{Interactive Visualization Design Tools}

Numerous interactive tools~\cite{Epigraphics, DataInk, Tableau, DataWrapper,
plotly, Lyra, DataIllustrator} have been developed to support visualization
authoring without conventional programming.

Chart typology tools and online chart makers \cite{manyeyes, DataWrapper,
rawgraphs, plotly} allow users to quickly create common charts from
data and compare design alternatives. More advanced systems such as
Tableau~\cite{Tableau} provide richer interactive authoring workflows,
including shelf-based encoding, automatic chart generation, and detailed
configuration of visual mappings. Nevertheless, these tools are primarily
oriented toward single-view visualizations or dashboard-style arrangements
and provide limited support for constructing composite visualizations through
explicit relationships among their constituent charts, such as layering,
faceting, and nesting.

Research systems further explore how direct manipulation and graphical editing
can support more expressive visualization design. Lyra~\cite{Lyra} combines
data-driven visualization specification with familiar vector-graphics-style
editing, while Data Illustrator~\cite{DataIllustrator} introduces lazy data
binding to flexibly transform graphical elements into data-driven graphics.
However, these systems primarily construct visualizations from low-level
graphical elements, such as rectangles, lines, and arcs, which enables highly
customized designs but may require users to construct and coordinate visual
structures from relatively primitive building blocks, even when the intended
design consists of recognizable charts.

In contrast, {\system} treats charts as the primary building blocks and
supports direct manipulation of chart-level composition relationships,
reducing the need to construct composite visualizations from low-level
graphical primitives.